\documentclass[11pt,a4paper]{article}
\usepackage[T1]{fontenc}
\usepackage[utf8]{inputenc}
\usepackage{amsmath,amssymb,graphicx,hyperref,longtable,booktabs,array}
\usepackage[margin=2.2cm]{geometry}
\providecommand{\tightlist}{\setlength{\itemsep}{0pt}\setlength{\parskip}{0pt}}
\setlength{\parindent}{0pt}\setlength{\parskip}{6pt}

\begin{document}
\hypertarget{piece-6-meal--and-insulin-driven-prior-26-nov-2026}{%
\section{Piece 6: meal- and insulin-driven prior (26 Nov
2026)}\label{piece-6-meal--and-insulin-driven-prior-26-nov-2026}}

Code: \texttt{code/inputs.py}, \texttt{code/test\_inputs.py} (9 of 9
pass), \texttt{code/run\_inputs.py}; results
\texttt{results/inputs\_results.json},
\texttt{results/inputs\_profiles.npz}. Written by Claude (AI assistant);
the team must read and sign off before claiming it.

\hypertarget{model}{%
\subsection{Model}\label{model}}

Blood glucose is driven by the logged meals and insulin instead of a
pure random-walk rate:

\[\frac{dG_I}{dt}=\frac{G_B-G_I}{\tau},\qquad \frac{dG_B}{dt}=R+a\,u_1(t)-b\,u_2(t),\qquad \frac{dR}{dt}=\sqrt{q}\,\xi(t),\qquad y_k=G_I(t_k)+e_k\]

\(u_1\) is carbohydrate (g/min) through a two-compartment gut kernel
(time constant 20 min); \(u_2\) is insulin above basal (U/min) through a
three-compartment subcutaneous kernel (time constant 40 min). The
patient gains \(a\) (mg/dL per g) and \(b\) (mg/dL per U) are extra
Kalman states with weak priors (5 and 40) and a slow random walk. Kernel
shapes are fixed population constants (my pick, not tuned). The 5-state
filter and smoother use an exact piecewise-constant discretisation of
the input term.

\hypertarget{tests-9-of-9-pass}{%
\subsection{Tests (9 of 9 pass)}\label{tests-9-of-9-pass}}

Zero inputs reproduce the Piece 4 filter (max difference 7e-14) and
smoother (7.5e-13); kernels normalised with peaks at 20 and 80 min;
basal insulin gives no drive and meal mass is conserved; on data
generated from the model the likelihood recovers tau (median 10.6 for
truth 10) and the gains (a 6.1 vs 6, b 48.3 vs 40); inputs cut smoother
RMSE (0.66 vs 3.25 mg/dL); smoother 95\% intervals cover 0.977 of the
truth; an AR(1) noise state raises the likelihood on AR(1)-noise data.

\hypertarget{results-27-simulated-patients-days-1-to-3-for-the-lag-days-4-to-5-for-evaluation}{%
\subsection{Results (27 simulated patients, days 1 to 3 for the lag,
days 4 to 5 for
evaluation)}\label{results-27-simulated-patients-days-1-to-3-for-the-lag-days-4-to-5-for-evaluation}}

Per-patient lag from the likelihood, truth = true effective lag:

\begin{longtable}[]{@{}llllll@{}}
\toprule\noalign{}
Condition & Prior & median abs. error & median signed error (min) & 95\%
interval covers truth & estimates at the 4-min floor \\
\midrule\noalign{}
\endhead
\bottomrule\noalign{}
\endlastfoot
S0, noise 2 & random walk & 50\% & -1.9 & 0.19 & 0.15 \\
& meal/insulin & 50\% & +4.0 & 0.22 & 0.04 \\
& placebo (inputs shifted) & 50\% & -2.5 & 0.19 & 0.22 \\
S0, noise 8 & random walk & 62\% & -6.4 & 0.19 & 0.78 \\
& meal/insulin & 51\% & -2.7 & 0.37 & 0.26 \\
& placebo & 63\% & -6.7 & 0.22 & 0.81 \\
S1, noise 2 & random walk & 52\% & -6.9 & 0.07 & 0.15 \\
& meal/insulin & 38\% & -1.0 & 0.33 & 0.04 \\
& placebo & 49\% & -7.9 & 0.11 & 0.19 \\
S1, noise 8 & random walk & 74\% & -11.4 & 0.11 & 0.78 \\
& meal/insulin & 66\% & -7.2 & 0.19 & 0.33 \\
& placebo & 75\% & -11.7 & 0.07 & 0.81 \\
\end{longtable}

Rapid-change RMSE (mg/dL, median over patients, days 4 to 5; paired
difference with 95\% CI for inputs minus random walk, tau = 10):

\begin{longtable}[]{@{}lllllll@{}}
\toprule\noalign{}
Condition & real-time filter, random walk & real-time filter, inputs &
difference & smoother, random walk & smoother, inputs & difference \\
\midrule\noalign{}
\endhead
\bottomrule\noalign{}
\endlastfoot
S0, noise 2 & 7.8 & 8.4 & +0.5 {[}-1.3, +1.4{]} & 4.9 & 4.6 & -0.3
{[}-0.3, -0.1{]} \\
S0, noise 8 & 20.3 & 14.9 & -5.0 {[}-8.5, -3.7{]} & 8.3 & 8.4 & -1.0
{[}-1.6, -0.2{]} \\
S1, noise 2 & 15.9 & 18.9 & +1.1 {[}+0.3, +4.7{]} & 13.5 & 11.9 & -1.2
{[}-1.3, -1.0{]} \\
S1, noise 8 & 25.8 & 21.2 & -5.0 {[}-8.6, -4.1{]} & 13.7 & 12.1 & -1.6
{[}-2.8, -1.0{]} \\
\end{longtable}

(Medians of differences need not equal differences of medians.)
Real-time hypoglycaemia detection is unchanged (23 of 23 events); false
alarms per day fall from 0.59, 1.04, 0.56, 1.17 to 0.26, 0.69, 0.22,
0.81 in the four conditions.

\hypertarget{what-this-means}{%
\subsection{What this means}\label{what-this-means}}

\begin{enumerate}
\def\labelenumi{\arabic{enumi}.}
\tightlist
\item
  The input prior is real information: moving the inputs by 400 samples
  (placebo) removes the gain, and with inputs the likelihood bias falls
  and the pile-up at the 4-min grid floor drops from 0.78 to 0.26 to
  0.33 at noise 8.
\item
  It does not fix the misspecification. Per-patient likelihood estimates
  still miss the true lag by 38 to 66\% (median), coverage of the 95\%
  interval stays at 0.19 to 0.37, and correlation with the truth is near
  zero. Using the per-patient estimate in the smoother is worse than tau
  = 10 in S0 (+5.7 {[}+0.3, +8.6{]} mg/dL at noise 2, +2.7 {[}+0.2,
  +4.2{]} at noise 8). Per-patient lag stays unidentifiable.
\item
  What it does buy is a better state estimate at a fixed tau: the
  smoother improves in all four conditions, and the real-time filter
  improves by about 5 mg/dL at noise sd 8 (the noise-limited regime of
  Piece 4), but not at noise 2.
\item
  Told the true effective lag, the input smoother reaches 2.3, 7.3, 3.6
  and 7.2 mg/dL: the headroom from knowing tau is still large, so the
  open problem is estimating tau, not the state prior.
\end{enumerate}

\hypertarget{caveats}{%
\subsection{Caveats}\label{caveats}}

Inputs are the simulator\textquotesingle s exact logged carbohydrate and
insulin; real logs are noisy and incomplete (missed boluses, estimated
carbs), so these gains are an upper bound. Kernel constants were chosen
once, not tuned on these results; Piece 8 checks sensitivity to them.
\end{document}
